% Originally: the \section{Solutions} of content/week-02.tex

\section{Solutions}
\label{sec:open_closed_solutions}

\begin{exercisesolution}[Proof of \cref{ex:ai_open_balls_are_open}]
% Generator: Claude Opus 5 (Medium)
Let $y \in B_r(x)$, so $d(y,x) < r$, and set $\varepsilon := r - d(y,x) > 0$. We claim that
$B_\varepsilon(y) \subseteq B_r(x)$, which is exactly what \cref{def:open_closed} requires.

Let $z \in B_\varepsilon(y)$, so $d(z,y) < \varepsilon$. By the triangle inequality
(\cref{def:metric_space}),
\[d(z,x) \leq d(z,y) + d(y,x) < \varepsilon + d(y,x) = \big(r-d(y,x)\big) + d(y,x) = r,\]
so $z \in B_r(x)$. Since $y \in B_r(x)$ was arbitrary, $B_r(x)$ is open.

\begin{remark}
The entire content is the choice $\varepsilon := r - d(y,x)$: the radius left over after
travelling from the centre out to $y$. Note that it shrinks to $0$ as $y$ approaches the edge of
the ball, so no \emph{single} $\varepsilon$ works for every $y$ at once. That is the same
\qt{the constant depends on the point} phenomenon which separates continuity from uniform
continuity in \cref{rem:continuity_vs_uniform_continuity}, met here for the first time.
\end{remark}
\end{exercisesolution}

\begin{exercisesolution}[Proof of \cref{ex:ai_disk_boundary}]
% Generator: Gemini 3.6 Flash (Medium)
Since $D$ is open, its interior is $\operatorname{int}(D) = D$. The closure of $D$ is $\overline{D} = \{(x,y) \in \mathbb{R}^2 \mid x^2 + y^2 \leq 1\}$. By \cref{def:interior_closure_boundary}, the boundary is:
\begin{align*}
\partial D &= \overline{D} \setminus \operatorname{int}(D) \\
&= \{(x,y) \in \mathbb{R}^2 \mid x^2 + y^2 \leq 1\} \setminus \{(x,y) \in \mathbb{R}^2 \mid x^2 + y^2 < 1\} \\
&= \{(x,y) \in \mathbb{R}^2 \mid x^2 + y^2 = 1\}.
\end{align*}
\end{exercisesolution}

% ===========================================================================
% Official sheet 2 (2.1--2.7). Derived independently, then checked against
% exercises/Sol2_Analysis2_eng.pdf.
% ===========================================================================

\begin{exercisesolution}[Solution to \cref{ex:2.2}]
% Generator: Claude Opus 5 (High)
\textbf{(a) and (c) are necessarily true; (b) and (d) are not.}

\textbf{(a) True.} Both inclusions come from general properties of the closure. Since
$Y_i \subseteq Y_1 \cup Y_2$ and the closure is monotone,
$\overline{Y_i} \subseteq \overline{Y_1\cup Y_2}$ for $i=1,2$, which gives \qt{$\supseteq$}.
Conversely $\overline{Y_1}\cup\overline{Y_2}$ is a \emph{finite} union of closed sets, hence
closed (\cref{prop:open_closed_operations}\textbf{(c)}), and it contains $Y_1\cup Y_2$; since
$\overline{Y_1\cup Y_2}$ is the smallest closed set containing $Y_1 \cup Y_2$
(\cref{def:interior_closure_boundary}), \qt{$\subseteq$} follows.

\textbf{(c) True,} and it is the easy half: $Y_1\cap Y_2 \subseteq Y_i$ gives
$\overline{Y_1\cap Y_2} \subseteq \overline{Y_i}$ for each $i$ by monotonicity, hence
$\overline{Y_1\cap Y_2} \subseteq \overline{Y_1}\cap\overline{Y_2}$.

\textbf{(b) False,} and therefore \textbf{(d) False} as well, (d) being (b) and (c) combined. In
$\mathbb{R}$ take
\[Y_1 := (0,1), \qquad Y_2 := (1,2).\]
Then $\overline{Y_1}\cap\overline{Y_2} = [0,1]\cap[1,2] = \{1\}$, whereas
$Y_1\cap Y_2 = \emptyset$ and so $\overline{Y_1\cap Y_2} = \emptyset$. The two sets meet only in
the limit, and the closure of the intersection cannot see a point belonging to neither set.

\begin{remark}[Where finiteness enters, again]
Part \textbf{(a)} is false for infinite unions, for exactly the reason given in
\cref{rem:finiteness_is_necessary}. Taking $Y_q := \{q\}$ for each $q \in \mathbb{Q}$,
\[\bigcup_{q\in\mathbb{Q}}\overline{Y_q} = \mathbb{Q}
  \qquad\text{but}\qquad
  \overline{\bigcup_{q\in\mathbb{Q}}Y_q} = \overline{\mathbb{Q}} = \mathbb{R}.\]
So the honest statement is: the closure commutes with \emph{finite} unions, and only
one-directionally with intersections. That is the same asymmetry that runs through the whole of
\cref{prop:open_closed_operations}.
\end{remark}
\end{exercisesolution}

\begin{exercisesolution}[Solution to \cref{ex:2.3}]
% Generator: Claude Opus 5 (High)
\textbf{(a), (b) and (c) are necessarily true; (d) is not.} It is efficient to prove
\textbf{(c)} first, since \textbf{(b)} then follows from it.

\textbf{(c) True.} Fix $x,y \in X$. For every $a \in A$ the triangle inequality gives
$d(x,a) \leq d(x,y) + d(y,a)$, and the left-hand side is bounded below by $d(x,A)$, so
\[d(x,A) \leq d(x,y) + d(y,a) \qquad\text{for every } a \in A.\]
Taking the infimum over $a$ on the right yields $d(x,A) \leq d(x,y) + d(y,A)$. Exchanging $x$
and $y$ and combining,
\[\big|d(x,A) - d(y,A)\big| \leq d(x,y),\]
so $d(\cdot,A)$ is \textbf{$1$-Lipschitz} (\cref{def:lipschitz_continuous}), in particular
continuous --- which is what the remaining parts use.

\textbf{(a) True.} If $A$ is closed then $A^c$ is open, so for $x \in A^c$ there is $r>0$ with
$B_r(x) \subseteq A^c$. No point of $A$ lies in that ball, i.e.\ $d(x,a) \geq r$ for every
$a \in A$, and taking the infimum gives $d(x,A) \geq r > 0$.

\textbf{(b) True.} By \textbf{(c)} the function $d(\cdot,A)$ is continuous, and
$M = d(\cdot,A)^{-1}\big([1,\infty)\big)$ is the preimage of a closed set under a continuous
map, hence closed (\cref{def:continuous}\textbf{(c)}, applied to complements).

\textbf{(d) False.} In $\mathbb{R}$ take $A := [0,1]\cup\{2\}$, whose interior
$A^\circ = (0,1)$ is non-empty as required. At the point $x := 2$,
\[d(x,A) = 0 \quad\text{(because } 2 \in A\text{)}, \qquad\qquad
  d(x,A^\circ) = \inf_{a\in(0,1)}|2-a| = 1.\]
The isolated point $2$ of $A$ contributes to $d(\cdot,A)$ but is invisible to $A^\circ$. The
hypothesis \qt{$A^\circ$ is non-empty} is far too weak, because it constrains $A$ nowhere near
$x$; the statement would become true under a hypothesis such as $A = \overline{A^\circ}$.
\end{exercisesolution}

\begin{exercisesolution}[Solution to \cref{ex:2.4}]
% Generator: Claude Opus 5 (High)
All with respect to the standard metric on $\mathbb{R}$, using
\cref{def:interior_closure_boundary}.
\begin{center}
\begin{tabular}{@{}c l l l l@{}}
& \textbf{$Y$} & \textbf{$\operatorname{int}(Y)$} & \textbf{$\overline{Y}$} & \textbf{$\partial Y$} \\
\hline
(1) & $[0,1]$ & $(0,1)$ & $[0,1]$ & $\{0,1\}$ \\
(2) & $\mathbb{Q}$ & $\emptyset$ & $\mathbb{R}$ & $\mathbb{R}$ \\
(3) & $\emptyset$ & $\emptyset$ & $\emptyset$ & $\emptyset$ \\
(4) & $(0,1)$ & $(0,1)$ & $[0,1]$ & $\{0,1\}$ \\
(5) & $[-1,1)\setminus\{0\}$ & $(-1,0)\cup(0,1)$ & $[-1,1]$ & $\{-1,0,1\}$ \\
(6) & $[0,\infty)$ & $(0,\infty)$ & $[0,\infty)$ & $\{0\}$ \\
(7) & $\{0\}$ & $\emptyset$ & $\{0\}$ & $\{0\}$ \\
(8) & $\left\{\tfrac1n : n \geq 1\right\}$ & $\emptyset$ &
  $\left\{\tfrac1n : n \geq 1\right\}\cup\{0\}$ &
  $\left\{\tfrac1n : n \geq 1\right\}\cup\{0\}$ \\
\end{tabular}
\end{center}

Three rows repay a second look.
\begin{itemize}
  \item \textbf{(1) versus (4).} The closed interval and the open interval have the same
    interior, the same closure and the same boundary. These three operations simply cannot
    distinguish a set from any other set squeezed between its interior and its closure.
  \item \textbf{(5).} Removing the single point $0$ from $[-1,1)$ adds $0$ to the boundary
    without changing the closure: a removed interior point becomes a boundary point. Note also
    that $1 \in \partial Y$ although $1 \notin Y$, while $-1 \in \partial Y$ and $-1 \in Y$ ---
    the boundary is indifferent to membership.
  \item \textbf{(2) and (8).} Both have empty interior, so $\partial Y = \overline{Y}$: once a
    set has no interior, the boundary swallows the entire closure. That is the general reason
    for $\partial\mathbb{Q} = \mathbb{R}$ in
    \cref{ex:interior_closure_boundary_more}\textbf{(c)}, and not an accident of $\mathbb{Q}$.
\end{itemize}
\end{exercisesolution}
